

\inb{
  Change letters: the $d$ in the model section refers to a size that is comparable to the code distance where here $d$ refers to the dimension of the Toric complex
} 

\section{The $(d,d^{\prime})$-Dimensional Toric Code and the SWEEP Rule}

\begin{definition}[$d$-Dimensional Toric Complex] 
  \label{def:toric_complex}
  For integers $d$ and $n$, denote the group resulting from the $d$-directed power of the cyclic group of order $n$ by $\mathcal{G}_{n}^{d} = \mathbb{Z}_{n}^{\times d}$, and denote the standard generating set of $\mathcal{G}_{n}^{d}$ by $S = \{ 1_{i} \colon i \in [d] \}$. We define the \textit{$d$-dimensional toric complex} to be the hypergraph obtained from the Cayley graph $\text{Cay}(\mathcal{G}_{n}^{d}, S)$ by taking its vertices as the $0$-faces and for any $i \in [d-1]$ we define the $i$-faces as follow: First, for a vertex $v \in \mathcal{G}_{n}^{d}$ and for a subset of generators $S^{\prime} \subset S$ of size $|S^{\prime}| = i$, the $i$-face rooted at $v$ and spanned by $S^{\prime}$, denoted by $\theta_{v, S^{\prime}}$, is:
  \begin{equation*}
    \begin{split}
      \theta_{v,S^{\prime}} \colonequals \bigcup_{I \subset S^{\prime}} \left\{  \left(\prod_{g \in I} g \right) v  \right\}
    \end{split}
  \end{equation*}
  The $i$-faces of the complex are the collection of all the possible rooted spanned $i$-faces induced by $\text{Cay}(\mathcal{G}_{n}^{d})$. We will denote them by $\mathcal{F}_{i}$. 

  Additionally, we name $n$ and $d$, the \textit{side length} and the \textit{dimension} of the complex.  
\end{definition}
\begin{definition}[Positive Edges and Cubes]
  For $u,v \in \mathcal{G}_{n}^{d}$, and $S^{\prime}_{1},S^{\prime}_{2} \subset [d]$, with sizes $|S^{\prime}_{1}| = |S^{\prime}_{2}| = i$, we say that the $i$-faces $\theta_{v,S^{\prime}_{1}}$ and $\theta_{u,S^{\prime}_{2}}$ are adjacent if the intersection $\theta_{v,S^{\prime}_{1}} \cap \theta_{u,S^{\prime}_{2}}$ is an $i-1$-face. In addition, if $\theta_{1}$ and $\theta_{2}$ are $i$ and $i-1$-faces such that $\theta_{2} \subset \theta_{1}$, then we say that $\theta_{2}$ is a \textit{positive edge relative to face} $\theta_{1}$ if they are rooted at the same vertex. Furthermore, we say that $\theta_{1}$ is a \textit{positive cube} of $\theta_{2}$. When the $i$-face is clear from the context, we will abbreviate and say that $\theta_{2}$ is \textit{positive}.
\end{definition}
\begin{example}
  For example, let $v, u \in \mathcal{G}_{n}^{d}$ and $g_{1} \neq g_{2} \in S$. Consider the 2-face $\theta = \{v, g_{1}v, g_{2}g_{1}v, g_{2}v\}$. Then, for $x, y \in \mathcal{G}_{n}^{d}$, we say that an edge (1-face) $\{x, y\}$ is \textit{positive relative to the face} $\theta$ if $x = v$ and $y$ is either $g_{1}v$ or $g_{2}v$.
\end{example}

We denote the complex by $G = (V, E, F)$, where $V$, $E$, and $F$ are the $0$, $1$, and $2$ faces.

\inb{Add a paragraph that explains that we choose not to use a chain-complex since we earn nothing from it i.e don't care about topology properties. Yet mention that chain-complex formulation proved itself as an effective tool to compute the properties of the codes we use. }

\begin{definition}[Assignment on The $(d^{\prime},d)$-Toric Complex]
  Let $\mathcal{T}$ be a $d$-dimensional complex. Each of its $i$-face sets naturally induces a vector space $\mathbb{F}_{2}^{|\mathcal{F}_{i}|}$. We call a binary vector $z \in \mathbb{F}_{2}^{|\mathcal{F}_{i}|}$ an assignment on the $i$-faces of $\mathcal{T}$. For a vertex $v$ and a subset of generators $S^{\prime} \subset S$ of size $|S^{\prime}|=i$, we use the notation $\rotfvs{}$ to represent the unit vector in $\mathbb{F}_{2}^{|\mathcal{F}_{i}|}$ that has $0$ in every coordinate except the one associated with the $i$-face $\rotfs{}$. Then
  \begin{equation*}
    \begin{split}
      \left\{ \rotfvs{} : v \in \mathcal{F}_{0}, S^{\prime} \subset S \text{ s.t } |S^{^\prime}| = i  \right\}
    \end{split}
  \end{equation*} is a base to $\mathbb{F}_{2}^{|\mathcal{F}_{i}|}$ \inb{ Not exactly base, they are linear dependent}.
Let $z$ be an assignment on the $i$-faces. Let $v_{1}, \ldots, v_{k} \in \mathcal{F}_{0}$ and $S_{1}, \ldots, S_{k} \subset S^{\times k}$ be the generator subsets such that $\theta_{v_{1},S^{\prime}_{1}}, \theta_{v_{2},S^{\prime}_{2}}, \ldots, \theta_{v_{k},S^{\prime}_{k}}$ are the collection of $i$-faces that support $z$, namely:
  \begin{equation*}
    \begin{split}
    z = \sum_{i=1}^{k}\rotfv{v_{i}}{S^{\prime}_{i}}
    \end{split}
  \end{equation*} 
  We refer to $\theta_{v_{1},S^{\prime}_{1}}, \theta_{v_{2},S^{\prime}_{2}}, \ldots, \theta_{v_{k},S^{\prime}_{k}}$ as the collection of faces that form $z$, or briefly, the faces form $z$.
  For $i > 0$, we define the linear map $\partial: \mathbb{F}_{2}^{|\mathcal{F}_{i}|} \rightarrow \mathbb{F}_{2}^{|\mathcal{F}_{i-1}|}$ such that
  \begin{equation*}
    \begin{split}
    \partial \rotfvs{} \colonequals \sum_{g \in S^{\prime}} \rotfv{v}{S^{\prime}/g} + \sum_{g \in S^{\prime}} \rotfv{gv}{S^{\prime}/g} 
    \end{split}
  \end{equation*}
  For $i < d$, we define the linear map $\partial^{*}: \mathbb{F}_{2}^{|\mathcal{F}_{i}|} \rightarrow \mathbb{F}_{2}^{|\mathcal{F}_{i+1}|}$ such that
  \begin{equation*}
    \begin{split}
      \partial^{*} \rotfvs{} \colonequals \sum_{g \in S/S^{\prime}} \rotfv{v}{S^{\prime} \cup g} + \sum_{g \in S^{\prime}} \rotfv{g^{-1}v}{S^{\prime}\cup g} 
    \end{split}
  \end{equation*}
  We define the \textit{restriction of $z$ to a vertex $v$} to be the linear map $|_{v} \colon \mathbb{F}_{2}^{|\mathcal{F}_{i}|} \rightarrow \mathbb{F}_{2}^{|\mathcal{F}_{i}|}$ such that: 
  \begin{equation*}
    \begin{split}
      \rotfv{u}{S^{\prime}} |_{v} \colonequals    
      \begin{cases}
      \rotfv{u}{S^{\prime}}  & v \in \theta_{u,S^{\prime}}  \\ 
        0 & \text{ else }
      \end{cases}
    \end{split}
  \end{equation*}
  We define the \textit{front of $z$ to a vertex $v$} to be the linear map $|_{v+} \colon \mathbb{F}_{2}^{|\mathcal{F}_{i}|} \rightarrow \mathbb{F}_{2}^{|\mathcal{F}_{i}|}$ such that: 
  \begin{equation*}
    \begin{split}
      \rotfv{u}{S^{\prime}} |_{v+} \colonequals    
      \begin{cases}
      \rotfv{u}{S^{\prime}}  & u = v \\ 
        0 & \text{ else }
      \end{cases}
    \end{split}
  \end{equation*}
  We say that an assignment $z$ is \textit{connected} if the graph that is induced by taking the faces that support $z$ as vertices and defining the edges to be the adjacency relation between them is connected. We define the connected components of $z$ to be its subsets such that each of them is a connected assignment. For a $j$-face $f$, we say that $f \in z$ if $z$ has a support on $\rotfvs{}$ such that $f \in \rotfs{}$. Consider an assignment $z \in \mathbb{F}_{2}^{\mathcal{F}_{i}}$ and a vertex $v \in \mathcal{G}_{n}^{d}$. We say that $z$ is a \textit{rooted assignment} at vertex $v$ if there are subsets $S^{\prime}_{1}, \ldots, S^{\prime}_{k} \subset S^{\times k}$ such that each of $S^{\prime}_{j}$ is at size $i$ and $z$ is decomposed as:
  \begin{equation*}
    \begin{split}
    z = \sum_{j}^{k}\rotfv{v}{S^{\prime}_{j}}
    \end{split}
  \end{equation*}
  Put differently, $z$ is the result of taking a collection of $i$-faces that are rooted at $v$.
\end{definition}

\inb{Add a proof for $\partial \rotfvs \cdot \partial^{*} \rotfv{u}{S^{\prime\prime}} = 0 \Rightarrow $  gives a CSS code.}

\inb{Need to decide on what should I elaborate, for example should I provide a proof for the distance and the rate of the code? }\cite{Dennis_2002}
\begin{definition}[$(d,d^{\prime})$ Toric Code] 
  Let $d, d^{\prime}$ be integers such that $1 < d^{\prime} < d$. For any integer $n$, we construct the $n$th instance of the quantum code family $(d,d^{\prime})$ \textit{Toric Code} as follows: Let $G = (\mathcal{F}_{0}, \ldots, \mathcal{F}_{d})$ be the $d$-dimensional toric complex obtained from $\mathcal{G}_{n}^{d}$ as defined in \Cref{def:toric_complex}. We associate the (q)bits with the $\mathcal{F}_{d^{\prime}}$, the $X$-checks with the $\mathcal{F}_{d^{\prime}-1}$, and the $Z$-checks with the $\mathcal{F}_{d^{\prime}+1}$. An $X$-check, $c_{x}$, is satisfied if the parity of (q)bits set on the $d^{\prime}$-face containing $c_{x}$ is even. Similarly, a $Z$-check, $c_{z}$, is satisfied if the parity, in the Hadamard basis, of the (q)bits set on the $d^{\prime}$-faces contained in $c_{z}$ is even.
\end{definition}

\begin{definition}[Sweep Rule]
  Let $z$ be an assignment over the $d^\prime$-faces, The \textit{sweep rule decoding procedure} at vertex $v \in \mathcal{G}_{n}^{d}$, defining as follow: 
  \begin{enumerate}
    \item For the $\mathcal{C}_{X}$ code. Look for a rooted assignment $\tilde{z}$ at $v$ such that $(\partial\tilde{z})|_{v+} = \partial(z)|_{v+}$. Namely, the $X$-checks rooted at $v$ get the same syndromes for $z$ and $\tilde{z}$. If there is one, then flip $\tilde{z}$. 
    \item Apply the Hadamard gate on any $d^{\prime}$ face containing $v$.
    \item For the $\mathcal{C}_{Z}$ code. Look for a rooted assignment $\tilde{z}$ at $v$ such that $(\partial^{*}\tilde{z})|_{v+} = \partial^{*}(z)|_{v+}$. Namely, the $Z$-checks rooted at $v$ get the same syndromes for $z$ and $\tilde{z}$. If there is one, then flip $\tilde{z}$.
    \item Apply the Hadamard gate again on any $d^{\prime}$ face containing $v$ to reverse the second stage.
  \end{enumerate}
  Stages $1$ and $3$ involve measurement of stabilizers.
\end{definition}

\begin{definition}[Decoding Gadget]
  Consider a quantum circuit $C$, for a quantum gate $u$ in $C$, with fun-in $\Delta$, we say that $u$ is a \textit{decoding gadget} if it has the following form: 
  \begin{enumerate}
    \item The wires of $u$ wires can be decomposed into two subsets $A$ and $B$ such that the wires in $A$ are reset wires, and there exists an encoded register $R$ in $C$ such that the wires in $B$ are all set in $R$.
    \item For a set of faults $\mathcal{E}$, there is a Pauli gate $u^\prime$ that acts only on $B$ such that the action of $u$ in $C[\mathcal{E}]$ equals $u^{\prime}$. Put differently, the quantum circuit obtained from $C[\mathcal{E}]$ by replacing $u$ with $u^{\prime}$ equals $C[\mathcal{E}]$.
  \end{enumerate}
\end{definition}


\begin{definition}[Unitary Decoding Representation]
Let $C = \{u_{i}\}$ be a quantum circuit, and let $\mathcal{U} = \{u^{\prime}_{i}\} \subset C$ be a subset of the gates in $C$ that are decoding gadgets. For a set of faults $\mathcal{E}$, we define the \textit{unitary decoding representation} of $C[\mathcal{E}]$ to be the quantum circuit obtained by replacing each of the gates in $\mathcal{U}$ with their corresponding Pauli gate according to $\mathcal{E}$.
\end{definition}


\begin{claim}
  The local gates performing Toom's rule are decoding gadgets. 
\end{claim}
\begin{proof}
  \inb{prove that.}
\end{proof}

\section{Infinite Torics}

\begin{definition}[Registers Coordinates System]
Let $C = \{u_{i}\}$ be a quantum circuit, with width $w$ and registers $\mathcal{R} = R_{1}, \ldots, R_{k}$, each of length at most $n$. For register $R_{i}$, we define the register identifier $\mathbf{1}_{R_{i}} = \mathbf{1}_{i}$. Then, we define the map $\phi_{\mathcal{R}} \colon \mathbb{Z}_{w} \rightarrow \mathbb{Z}_{R} \times \mathbb{Z}_{n}$ to return, for a space-location coordinate $x \in \mathbb{Z}_{w}$, the tuple consisting of the identifier of the register containing $x$, and its relative coordinate within the register. Namely,
\begin{equation*}
  \begin{split}
    \psi_{\mathcal{R}}(x) = (\mathbf{1}_{R_{j}}, x^{\prime})
  \end{split}
\end{equation*} 
Where $R_{j}$ is the register containing $x$, and $x^{\prime}$ is the index of the qubit, set at the $x$ location, in $R_{j}$.
The presentation of $C$ in the \textit{registers coordinate system} is the encoding by a set of gates $\{u^{\prime}_{j}\}$, which is the result of replacing each space-location of the original gates with the applications of $\psi_{\mathcal{R}}$ on them. Namely, for each $u_{i} = (g_{i}, l_{i}) \in C$, we map its space-location coordinate $l_{i_{2}} = (l_{i_{2}}^{(1)}, \ldots, l_{i_{2}}^{(\Delta)})$ to:
      \begin{equation*}
        \begin{split}
          l_{i_{2}} \mapsto ( \psi_{\mathcal{R}}(l_{i_{2}}^{(1)}), ..,\psi_{\mathcal{R}}( l_{i_{2}}^{(\Delta)}))
        \end{split}
      \end{equation*}
For brevity, we name the first and second coordinates in the result of $\psi_{R}$ the \textit{register coordinate} and the \textit{inner coordinate}.
    \end{definition}
    \begin{definition}[Shifting in Registers System]
      
Consider the register coordinate system induced by a quantum circuit of width $w$, with $R$ registers, each of length at most $n$. For an integer $r \in \mathbb{Z}$, we define the $r$-shifting function $\text{Sh}_{r} \colon (\mathbb{Z}_{R} \times \mathbb{Z}_{n})^{\times \Delta} \rightarrow (\mathbb{Z}_{R} \times \mathbb{Z})^{\times \Delta}$ to act on space-location coordinates as follows:
  \begin{equation*}
    \begin{split}
      & \text{Sh}_{r}( l_{1}, .., l_{k} ) = ( \text{Sh}_{r}(l_{1}), .., \text{Sh}_{r}(l_{k}) ) \\ 
      & \text{ Where } \text{Sh}_{r}(l_{i}) = 
      \begin{cases}
        l_{i} + r & \text{ if } l_{i} \text{ is an inner coordinate } \\ 
        l_{i}& \text{ else }
      \end{cases}
    \end{split}
  \end{equation*}
\end{definition}

\begin{definition}
  \label{def:extension}
Let $n$ be an integer, and let $C = \{u_{i}\}$ be a quantum circuit with registers $\mathcal{R} = R_{1}, \ldots, R_{k}$ such that each register $R_{i}$ is a toric encoding with a side length $n$. In addition, all the decoding gadgets of $C$ are sweep rule gadgets; let us denote them by $\mathcal{U}$. We define the \textit{$\infty$-extension of $C$}, denoted by $C^{\infty}$, as follows:
  \begin{enumerate}
    \item For a register $R_{i} \in \mathcal{R}$, denote by $d_{i}, d_{i}^{\prime}$ the dimension of the toric complex that induces the code, and the dimension of the faces that are associated with the bits. Each register $R_{i}$ in $\mathcal{R}$ is mapped to a register $R^{\prime}_{i}$, which is a $(d_{i}, d_{i}^\prime)$-dimensional, infinite side length, toric encoding. Since the mapping between registers is one-to-one, we use the identifier $\mathbf{1}_{R_{i}}$ to identify also $R^{\prime}_{i}$.
    \item Each gate $u_{i} \in C / \mathcal{U}$ is mapped to:     
      \begin{equation*}
        \begin{split}
          u_{i} & = (g_{i},l_{i}) \mapsto  
          \bigcup_{j = -\infty}^{\infty}\left\{ (g_{i},( \text{Sh}_{jn}(l_{i})) \right\}
        \end{split}
      \end{equation*}
    \item Each gate $ u \in \mathcal{U} $ is mapped to the collection of sweep rule gadgets $\{ v_{j} \}_{-\infty}^{\infty}$ such that each of them has the same time coordinate as $ u $, is in the register that corresponds to the image of $ u $'s register, and is rooted at $\text{root}(v) + j n$.
  \end{enumerate}
For a set of faults $\mathcal{E} = \{ f_{i} \}$, we denote by $C^{\infty}[\mathcal{E}^{\infty}]$ the \textit{$\infty$-extension of $C[\mathcal{E}]$}.
\end{definition}


\begin{claim}
\label{claim:pauli_back}
  Let $ u = (g,l)$ and $v_{0},v_{\pm 1}$ plantation of $u$ at locations $\text{root}(u)$ and $\text{root(u)} \pm n$. Then $v_{0}v_{\pm 1}|_{[w]} = u$.
\end{claim}
\begin{proof}
  \inb{Prove it}
\end{proof}

\begin{definition}
  For a quantum circuit $C$, of the form defined in \Cref{def:extension}, and a set of faults $\mathcal{E}$, denote:
  \begin{equation*}
    \begin{split}
      \rho & \colonequals C[\mathcal{E}] \ketbra{0}{0} C[\mathcal{E}]^{\dagger} \\
      \rho_{\infty} & \colonequals C^{\infty}[\mathcal{E}^{\infty}]\ketbra{0}{0}\left(C^{\infty}[\mathcal{E}^{\infty}]\right)^{\dagger}
    \end{split}
  \end{equation*}
We say that a quantum circuit $C$ is \textit{$\infty$-compatible} if, for any set of faults $\mathcal{E} = \{ f_{i} \}$, it holds that:
  \begin{equation*}
    \begin{split}
      \rho =  \Trr{}_{/[w]}\left( \rho_{\infty} \right) 
    \end{split}
  \end{equation*}
\end{definition}
\begin{claim}
  Let $C$ be quantum circuit of the form defined in \Cref{def:extension}. Then $C$ is $\infty$-compatible.
\end{claim}
\begin{proof}
  Its enough to prove the claim for circuits in their unitary decoding representation. We prove it by induction on the number of the gates.

  \begin{enumerate}
    \item Base. For an empty circuit, $\rho_{\infty}$ is just the infinite string of zeros; thus, the restriction of it to $[w]$ locations, namely $ \Trr{}_{/[w]}( \rho_{\infty} ) $, yields a single matrix element corresponds to the $w$-length string of zeros, which is exactly $\rho$.
    \item Step. Assume the correction of the claim for any circuit, with less than $\tau$ steps, and consider a circuit $C$ with $\tau$ steps. Denote by $C^{\prime}$ the circuit obtained form $C$ by erasing the last gate in $C$, observes that $C^{\prime}$ is a circuit with $\tau - 1$ steps. Denote by $\rho^{\prime}$ and $\rho^{\prime}_{\infty}$ the: \inb{Enter it}.  
      By the induction assumption we have:
  \begin{equation*}
    \begin{split}
      \rho^{\prime} =  \Trr{}_{/[w]}\left( \rho^{\prime}_{\infty} \right) 
    \end{split}
  \end{equation*}
Denote the last gate in $C$ by $u$, and by $I = \{v_{j}\}_{-\infty}^{\infty}$ the collection of gates in $C^{\infty}$ to which $u$ is mapped. Let us split into cases, depending on whether $u$ is or is not a decoding gate.
  \begin{enumerate}
    \item Suppose that $u$ is not a decoding gate. Let $g$ and $l$ be the gate element and the location of $u$. By the definition of $\infty$-extension, $I$ equals:
      \begin{equation*}
        \begin{split}
              I = \bigcup_{j = -\infty}^{\infty}\left\{ (g,( \text{Sh}_{jn}(l)) \right\}
        \end{split}
      \end{equation*}
      So there is only one gate in $I$, corresponding to $j=0$, that has support on qubits with spatial location in the range $[w]$. Denote that gate by $v_{0}$. 

 \begin{equation*}
   \begin{split}
   \Trr{}_{/[w]}\left( \rho_{\infty} \right) &= \Trr{}_{/[w]}\left( \prod_{j = -\infty}^{\infty} v_{j} \rho_{\infty}^{\prime} \prod_{j = -\infty}^{\infty} v_{j}^{\dagger} \right) \\ 
   &= v_{0}\Trr{}_{/[w]}\left( \prod_{j = -\infty, j \neq 0}^{\infty} v_{j} \rho_{\infty}^{\prime} \prod_{j = -\infty, j \neq 0}^{\infty} v_{j}^{\dagger} \right) v_{0}^{\dagger} \\
   &= v_{0}\Trr{}_{/[w]}\left( \prod_{j = -\infty, j \neq 0}^{\infty} v_{j}^{\dagger} \prod_{j = -\infty, j \neq 0}^{\infty} v_{j} \rho_{\infty}^{\prime} \right) v_{0}^{\dagger} \\
   &= v_{0}\Trr{}_{/[w]}\left( \rho_{\infty}^{\prime} \right) v_{0}^{\dagger} \\
   &= v_{0} \rho^{\prime} v_{0}^{\dagger} \\
   &= \rho 
   \end{split}
 \end{equation*}

 \item In the case that $u$ is a decoding gate, observe that there are at most two gates in $I$ that have support on the qubits with space-location in $[w]$. If there is exactly one, then repeating the non-decoding gate case while setting this gate as $v_{0}$ gives the desired result. 
Else, in the case where two gates act non-trivially on qubits at $[w]$, denote those gates by $v_{0}, v_{1}$. Since we assumed that $C$ is given in its unitary decoding representation, it holds that $v_{0}, v_{1}$ are Paulis, in particular a product operator. Denote their product decompositions by $v_{0} = \prod_{k}v_{0}^{(k)}$ and $v_{1} = \prod_{k}v_{1}^{(k)}$. Thus, by repeating the non-decoding case, but extracting from the infinite product only the terms in $v_{0}$ and $v_{1}$ that act non-trivially at $[w]$, we get:
   \begin{equation*}
     \begin{split}
       \Trr{}_{/[w]}\left( \rho_{\infty} \right) &= v_{0}v_{1}|_{[w]} \rho^{\prime} \left(v_{0}v_{1}|_{[w]}\right)^{\dagger} 
     \end{split}
   \end{equation*}
Moreover, since $v_{0}$ and $v_{1}$ are a plantation of $g$ in locations $\text{root}(u)$ and $\text{root}(u) \pm n$, then by \Cref{claim:pauli_back}, we have that $u = v_{0}v_{1}|_{[w]}$; therefore:
   \begin{equation*}
     \begin{split}
       \Trr{}_{/[w]}\left( \rho_{\infty} \right) &= v_{0}v_{1}|_{[w]} \rho^{\prime} \left(v_{0}v_{1}|_{[w]}\right)^{\dagger} \\
       &= u \rho^{\prime} u^{\dagger} \\
   &= \rho 
     \end{split}
   \end{equation*}
  \end{enumerate}
  \end{enumerate}
\end{proof}

\inb{ add that the computation graph is lifted naturally to the $\infty$-extension, since the image of a mapping are overlapped gates, we can execute them simultaneously.  }  

\section{Shrinking}


\inb{ consider changing the name to something like hight lines / contour in order to hint on the similarity to topography map. In addition, the embedding inside $\mathbb{R}^{d}$ should appears before that definition. }
\begin{definition}[Potential Walls]
  We define the potential walls to be a functions collection from the vertices of $\mathcal{G}_{n}^{d}$ to the reals as follow:
  \begin{enumerate}
    \item For any coordinate $i \in [d]$, we define $L_{i} \colon \mathcal{G}_{n}^{d} \rightarrow \mathbb{R}$ to act as $L_{i}(v) = v_{i}$.
    \item For any subset of coordinates $I \subset [d]$ we define $L_{I} \colon  \mathcal{G}_{n}^{d} \rightarrow \mathbb{R}$ to act as $L_{I}(v) = -\sum_{i \in I}v_{i}$. 
  \end{enumerate}
For a subset of vertices $\Gamma \subset \mathcal{G}_{n}^{d}$, a vertex $v \in \Gamma$, and a potential wall $l \colon \mathcal{G}_{n}^{d} \rightarrow \mathbb{R}$, we say that $v$ is a \textit{local maximum} of $l$ in $\Gamma$ if for any neighbor $u$ of $v$, such that $u \in \Gamma$, we have $l(v) \ge l(u)$. We say that $v$ is a strong local maximum if the inequality is strict, namely $l(v) > l(u)$. When the subset $\Gamma$ is clear from the context, we might omit it and briefly say that $v$ is a local maximum of $l$.

\end{definition}

\inb{
  What we want is that if a bit contains both $v, gv$ then there will be a check that contains both $v, gv$, that is true if $d^{\prime}  > 1$, then it means that we can earse from the bit a vertex which it is not $gv$. 

}



\begin{claim}
  \label{claim:localmax}
Let $z$ be an assignment, and let $\Gamma$ be the set of vertices adjacent to unsatisfied checks induced by $z$. Let $v$ be a vertex in $\Gamma$, and let $g$ be a coordinate in $S$ such that $v$ is a local maximum of $L_{g}$. Then $gv \notin \Gamma$.
\end{claim}
\begin{proof}
  Suppose $gv \in \Gamma$, then 
  \begin{equation*}
    \begin{split}
      L_{g}(gv) = (gv)_{g} = 1 + v_{g} = 1 + L_{g}(v)
    \end{split}
  \end{equation*}
In contradiction to $v$ being a local maximum of $L_{g}$ in $\Gamma$.
\end{proof}
\begin{claim}
Let $z$ be an assignment, and let $\Gamma$ be the set of vertices adjacent to unsatisfied checks induced by $z$. Let $v$ be a vertex in $\Gamma$, and let $g$ be a coordinate in $S$ such that $v$ is a local maximum of $L_{g}$. Then for any check associated with a $d^{\prime}-1$ face rooted at $v$ that contains $gv$, it is satisfied.
\end{claim}
\begin{proof}
Assume, through contradiction, the existence of an unsatisfied check associated with a $d^{\prime}-1$ positive face rooted at $v$, which contains $gv$. This implies $gv \in \Gamma$, contradicting \Cref{claim:localmax}.
\end{proof}


\begin{claim}
  Let $z$ be an assignment. There is $\tilde{z} \sim_{\mathcal{C}_{X}} z$ such that for any vertex $v \in \tilde{z}$, which is a local maximum of $L_{d+1}$ in $\Gamma_{\tilde{z}}$, it holds that $v$ is also a local maximum of $L_{d+1}$ in $\tilde{z}$.
\end{claim}
\begin{proof}
  \inb { For $d = 4, d^{\prime}= 2$.} Denote the bits, rooted at $v$ and spanned by generators $g_{i},g_{j}$ by $g_{ij} = g_{ji}$. 
  The restrictions: 
  \begin{equation*}
    \begin{split}
      r_{i} \colon \sum_{j\neq i} g_{ij}
    \end{split}
  \end{equation*}
  Thus we get
  \begin{equation*}
    \begin{split}
      \sum_{i}^{n}r_{i} & = \sum_{i}^{n} \sum_{j \neq i }^{n+1} g_{ij} \\
      & =   \sum_{i}^{n}\sum_{j < i}^{n} g_{i,j} + g_{j,i} +   \sum_{i}^{n}  g_{i,n+1} \\
      & = r_{n+1}
    \end{split}
  \end{equation*}
  Thus the dimension of the small code at $v$ is at least: 
  \begin{equation*}
    \begin{split}
      { d \choose 2 } - (d-1) = 3 
    \end{split}
  \end{equation*}
  We have $ { 4 \choose 3 } = 4 $ stabilizers rooted at $v$ (each for each cube). Yet:  
  \begin{equation*}
    \begin{split}
      s_{1} &= g_{1,2} + g_{2,3} + g_{1,3} \\
      s_{2} &= g_{1,2} + g_{2,4} + g_{2,4} \\
      s_{3} &= g_{1,3} + g_{3,4} + g_{1,4} \\
      s_{4} &= g_{2,3} + g_{3,4} + g_{2,4} 
    \end{split}
  \end{equation*}
  And clearly 
  \begin{equation*}
    \begin{split}
      s_{4} = s_{1}+s_{2}+s_{3}
    \end{split}
  \end{equation*}
\end{proof}




\begin{claim}
  Let $z$ an assignment rooted at vertex $v$, such $\partial z |_{v+} = 0$. Then there is a stabilizer $w$, rooted at $v$, such that $w|_{v+} = z$.  
\end{claim}

\begin{proof}
Denote by $F$ the faces in $z$, and by $J$ the collection of generators that support $F$. We claim that $|J| \ge d^{\prime}+1$.

Clearly, $|J| \ge d^{\prime}$, since each face is defined by a root vertex and $d^{\prime}$ generators. Yet, observe that if $|J|$ equals exactly $d^{\prime}$, then $z$ contains a single $d^{\prime}$-face, denote it by $f$. Then, any check associated with a $d^{\prime}-1$ face of $f$ is unsatisfied, in particular the $d^{\prime}-1$ faces that are rooted at $v$, in contradiction to $\partial z|_{v+} = 0$. 


Now, one can choose $J^{\prime} \subset J$ of size $d^{\prime}+1$, and define the stabilizer $w = \theta_{v,J^{\prime}}$.
    
\end{proof}


\begin{claim} \label{claim:syn_edge}
Suppose $v \in z$ and for any $g \in S$, we have $g^{-1}v \notin z$, then there is $\rotf{v}{J} \in \partial z$.
\end{claim}
\begin{proof}
  
  \begin{equation*}
    \begin{split}
      \partial z = \sum_{S^{\prime}} \rotfv{v}{S^{\prime}} + .. 
    \end{split}
  \end{equation*}
\end{proof}


\begin{claim}
  Let $z$ be a connected assignment on the $d^{\prime}$-faces at finite size, and let $\Gamma$ be the set of vertices which are adjacent to unsatisfied checks induced by $z$. Let $v$ be a vertex in $\mathcal{G}^{d}_{\infty}$ such that $v$ is a strong global maximum of $L_{d+1}$, at $\Gamma$. Then there is a rooted assignment $\tilde{z}$ at $v$ such that $z_{v} = \tilde{z}_{v}$. 
\end{claim}
\begin{proof}
  Let $ \Theta = \theta_{v_{1},S^{\prime}_{1}}, \theta_{v_{2},S^{\prime}_{2}},.. \theta_{v_{k},S^{\prime}_{k}}$ be the faces forms $z$. 


  Since $z$ is finite, it has a maximum for $L_{d+1}$. Denote by $u$ the vertex that gets the maximum, and denote by $\theta_{u,J_{1}}, \ldots, \theta_{u,J_{l}}$ the $d^{\prime}$-faces in $z$ that are rooted at $u$. We claim that $u = v$.

  Suppose not, namely $u \notin \Gamma$. Then it means that any $d^{\prime}-1$ face containing $u$ is satisfied. Observe that $\sum \theta_{u,J_{i}} \neq 0$; otherwise, it would contradict the non linearity independence of  $\rotfv{v_{1}}{S^{\prime}_{1}}, \rotfv{v_{2}}{S^{\prime}_{2}},.. \rotfv{v_{k}}{S^{\prime}_{k}}$. So, by \Cref{claim:syn_edge} there exists a $d^{\prime}-1$ face rooted at $u$, contained in $\sum \theta_{u,J_{i}}$. Let us denote it by $\theta_{u,J^{\prime}}$.

Since $\theta_{u,J^{\prime}}$ is satisfied, it has to be adjacent to a turned bit, namely a $d^{\prime}$ face in $z$, not in $\theta_{u,J_{1}}, \ldots, \theta_{u,J_{l}}$. Denote that $d^{\prime}$ face by $f$. Yet, since $\theta_{u,J^{\prime}}$ is rooted at $u$, it contains $u$, and therefore $f$ contains $u$. Now, if $f$ is rooted at $u$, then it is a contradiction that $\theta_{u,J_{1}}, \ldots, \theta_{u,J_{l}}$ are all the $d^{\prime}$ faces in $z$ rooted at $u$. If $f$ is not rooted at $u$, denote its root by $w$, so there is $g \in S$ such that $gw = u$, and $w \in z$, in contradiction to $u$ being maximum for $L_{d+1}$.


So, we find that $v$ is also the maximum of $L_{d+1}$ in $z$. Therefore, we have that $\tilde{z} = \sum \theta_{u, J_{i}}$ zeros out $z_{v}$.
\end{proof}



\newpage



\begin{claim}
  Let $z$ be a connected assignment on the $d^{\prime}$-faces at finite size, and let the collection $\Theta = \theta_{v_{1},S^{\prime}_{1}}, \theta_{v_{2},S^{\prime}_{2}},.. \theta_{v_{k},S^{\prime}_{k}}$ be the decomposition forms $z$. Denote by $\Gamma$ the set of vertices adjacent to unsatisfied check induced by $z$. Then $\Theta$ is acyclic, if and only if $z$ has at least one Infimum. 
\end{claim}
\begin{proof}
  First, suppose that $z$ has at least one Infimum. Let $u$ be an infimum in $z$.. 
  Second, suppose that $\Theta$ is acyclic. 
\end{proof}



\begin{claim}\label{claim:extermal}
Suppose that $d^{\prime} > 1$ and a vertex $v \in \mathcal{G}^{d}_{\infty}$ is also adjacent to the domain wall, namely $v \in \Gamma$. It holds that $v$ cannot be simultaneously both a (strong) local max of $L_{d+1}$ and a local min of $L_{d+1}$.
\end{claim}

\begin{proof}
  Let $v \in \Gamma$ be a local maximum point of $L_{d+1}$, It means that:
  \begin{enumerate}
    \item The vertex $v$ is adjacent to a $d^{\prime - 1}$-face which is unsatisfied. 
   \item That face is rooted in $v$; otherwise, the root of the face is a vertex in $\Gamma$ with a higher value of $L_{d+1}$, in contradiction to $v$ being a local maximum.
  \end{enumerate}
Thus, there is a subset of generators $S^{\prime} \subset S$, of size $d^{\prime} - 1 > 0$, such that the unsatisfied check set is on the face $\theta_{v,S^{\prime}}$. Given $d^{\prime} > 1$, we have that $S^{\prime}$ contains at least one more element besides the empty set, denoted by $g$. Now, consider the vertex $gv$. Since $gv$ is adjacent to $\theta_{v,S^{\prime}}$, then $gv \in \Gamma$. On the other hand, on the infinite toric $\mathcal{G}^{d}_{\infty}$, moving along the grid generator decreases $L_{d+1}$ by one, so $L_{d+1}(gv) = L_{d+1}(v) - 1$. In other words, $gv$ is a neighbor of $v$ with a lower value of $L_{d+1}$, and thus $v$ cannot be a local minimum of $L_{d+1}$.
\end{proof}

\begin{remark}

\inb{
  That claim alone gives a decoder in the clean setting. Write a remark indicating that when $d^\prime = 1$, meaning qubits are on the edges (2D toric), the claim is not true.
}
\end{remark}


\begin{definition}[Reduced subset with boundrey at $v$]
  
\end{definition}
